\chapter{Institutions: The Room for the River Case}
\label{ch:room-for-river}

\begin{plainlanguage}
\textbf{In plain terms:} Does a multi-part system (here, Dutch flood infrastructure) adapt to crisis all at once, or at different speeds for different parts? Closely linked parts changed almost instantly; the surrounding institutions turned out to be on their own centuries-long timeline unrelated to the crisis. The chapter also catches systems that \emph{announce} change without actually changing underneath.
\end{plainlanguage}

\begin{chapternote}
A note on what Part~III is for. A reader could reasonably ask why a book about defining intelligence is about to spend two chapters on Dutch flood infrastructure and the way small children mislabel cows as dogs. The answer is that Part~II built a piece of machinery --- fidelity-weighted elimination, converging toward a hard core and a soft periphery --- and machinery built in the abstract earns nothing until it has been pointed at something messy and watched to see what happens. These two chapters are not detours on the way to the real argument. They are the moment the book finds out whether its own formalism survives contact with cases it was not designed around. The first naive prediction tested here fails, instructively, in both domains. What survives the failure is more useful than what the naive version would have given if it had simply been confirmed.
\end{chapternote}

\section{A clean prediction, stated in advance}

Take a system with multiple coupled subsystems undergoing repair after a shared triggering failure, and ask: do they reach closure --- do they stabilize a new, working response to the failure --- at the same time, or at different times? If different times, is there a predictable pattern to the gap?

A natural first hypothesis, building directly on the active-inference-style alternative the fidelity framework was built to out-predict, would be that subsystems closer to the proximate cause of failure close first, and subsystems further from it --- more socially or institutionally mediated --- close later. Applied to a flood-control system with hydraulic, ecological, and institutional/governance subsystems, this predicts an ordering:
\[
t_H < t_E < t_I
\]
hydraulic recognition before ecological recognition before institutional reorganization, on the reasoning that the physical failure is most immediate and the governance response most mediated and slowest.

This prediction is stated in advance, plainly, because what happens to it next is the chapter's actual content.

\section{The case}

The Netherlands' Room for the River program offers an unusually well-documented test. Severe floods on the Rhine and Meuse in 1993 and 1995 forced the evacuation of roughly a quarter million people, and the floods are widely credited as the starting point of a fundamental policy reconsideration --- away from the centuries-old strategy of progressively raising and strengthening river embankments, toward a strategy of deliberately surrendering floodplain land back to the rivers. The full program, formally adopted via a 2006 Spatial Planning Key Decision, ran construction from 2007 to roughly 2015--2018 across more than thirty project sites.

The immediate post-flood response was not the eventual program. It was incremental: an adjustment of the design discharge standard the embankment system was built to handle, from 15{,}000 to 16{,}000 cubic meters per second --- a parameter change within the existing engineering framework, not yet a structural reorganization. The decade between flood and formal structural decision is itself worth treating as data: it is roughly the gap between the triggering failure and the point at which the underlying engineering logic (``keep raising the embankments'') was recognized, not as having failed once, but as having no terminus --- a closure-exhaustion realization, arrived at only after the incremental route had visibly been tried and had visibly not resolved the deeper problem.

\section{Where the ordering hypothesis breaks}

The predicted ordering $t_H < t_E < t_I$ does not survive contact with the documentary record, and it breaks in two distinct and instructive ways.

First: hydraulic and ecological closure were not sequential at all. The program was explicitly designed, from early in its formation, around a joint objective balancing hydraulic effectiveness, ecological robustness, and spatial-cultural quality together --- not ecological closure trailing years behind hydraulic closure, but the two being defined as a single combined admissibility criterion from the outset.

\begin{definition}[Interface Lag]
For two coupled subsystems $S_i, S_j$ with closure times $t_i$ and $t_j$, define $\Delta t_{ij} = |t_i - t_j|$.
\end{definition}

\subsection*{The ontology of subsystems}

Before going further it is worth being explicit about what a ``subsystem'' is, since the word is about to carry real argumentative weight, here and even more so in Chapter~\ref{ch:development}, where subsystem-indexed repair becomes the chapter's central repair move. Three questions are worth separating. Are subsystems discovered, existing in the world independently of anyone's analysis, or are they analyst-imposed, a partition the observer brings to the case rather than finds in it? Can subsystem boundaries change over time, or are they fixed once drawn? And what stops ``subsystem'' from becoming an infinitely flexible category, redrawn however convenient whenever a prediction needs rescuing?

This book's working position is deliberately modest rather than metaphysically ambitious: a subsystem is whatever a perspective (in the Chapter~\ref{ch:candidate-space} sense) can independently measure and report on. Hydraulic engineering and institutional governance are treated as separate subsystems in this chapter not because reality is carved at uncontroversial joints, but because each has its own independent measurement methodology, its own historical record, and its own closure criterion that can be assessed without reference to the other --- which is precisely the anti-vacuity standard this book holds subsystem-splitting to throughout, stated here first and then enforced explicitly as a named Principle when Chapter~\ref{ch:development} needs it. On this view subsystem boundaries are not fixed for all time; they can shift as measurement capability shifts, which is itself informative rather than a defect --- a boundary that only becomes visible once a new measurement technique exists is not arbitrary, it is discovered late. What rules out the worry that ``subsystem'' becomes infinitely flexible is exactly the discipline a reader should expect by now: no subsystem split is legitimate in this book without an independently motivated measurement method behind it, and chapters that introduce a split are expected to show that method, not merely assert the split.

\begin{conjecture}[Coupling--Lag Relation]
$\Delta t_{ij}$ is inversely related to the coupling strength $C_{ij}$ between the subsystems --- tightly coupled subsystems should reach closure near-simultaneously, while loosely coupled ones should drift apart.
\end{conjecture}

The hydraulic--ecological pair in this case shows near-zero $\Delta t_{HE}$, consistent with high coupling --- but this is not independent confirmation of the conjecture, since the two subsystems were defined together by program design rather than discovered to have converged. What the case actually offers here is an illustration of the high-coupling, low-lag corner of the conjecture, not a test of it: there was no opportunity for these two tracks to drift apart and then resynchronize, because they were never separated to begin with.

Second, and more substantively informative: the institutional/governance track does not share the 1993--1995 trigger at all. Dutch regional water-governance institutions --- the canoe-house-era water boards, their consolidation from roughly 3{,}500 around 1900 to about 2{,}600 by 1950 to 21 today --- trace a centuries-long reorganization driven by the growing power of national water-management authority since the late eighteenth century, a process already mid-stream when the 1990s floods occurred. Room for the River began as an ad hoc engineering reaction to a specific flood event and only afterward grew into application of a broader river-basin governance approach --- meaning the institutional reorganization here was not delayed-response-to-a-shared-trigger at all. It was an independent, much older process the flood event happened to intersect and accelerate, not originate.

This breaks an assumption buried in the original hypothesis: that all subsystems share a common reference event from which a lag can even be measured. Where they don't, $\Delta t_{ij}$ as defined is not well-formed --- there is no single $t=0$ from which both $t_H$ and $t_I$ can be measured, because the institutional clock was already running on its own schedule before the hydraulic crisis began.

\section{A second, separate finding: nominal versus operational closure}

A third phenomenon surfaces in the record, distinct from lag and worth a definition of its own. Official Dutch water policy doctrine proclaims water as the ordering principle for spatial planning --- a stated, adopted closure at the level of policy language --- yet actual spatial-development decisions are reported as continuing to follow prior economic and institutional priorities, with researchers concluding the resulting centralization tendency runs counter to what the new principles intended. The principle was formally adopted; the underlying practice it was meant to govern did not actually reorganize beneath it.

\begin{definition}[Nominal vs.\ Operational Closure]
Let $N(t) \in \{0,1\}$ denote whether a closure has been formally declared or adopted in a system at time $t$, and $O(t) \in \{0,1\}$ denote whether the underlying practice has actually reorganized to match it. A \emph{closure illusion} is a state in which $N(t) = 1$ while $O(t) = 0$ --- the system has said it closed without having closed.
\end{definition}

This is a different failure mode than lag. Lag is a question of timing between subsystems that do eventually both close. Closure illusion is a question of whether a declared closure is real at all, and it requires measuring $N$ and $O$ independently rather than treating a policy announcement as evidence of operational change --- a methodological caution worth carrying into every later chapter that draws on institutional, cultural, or self-reported evidence of closure, since self-report is exactly where $N$ and $O$ are most easily conflated.

A related complication is worth naming rather than setting aside: several independent sources on flood infrastructure converge on the observation that structural hardening measures (levees, dams) tend to suppress the very signal that would otherwise prompt structural reorganization, because the public perceives risk as resolved and development on the protected floodplain increases as a result. Call this \emph{patch-induced closure illusion} --- a local fix that does not merely fail to trigger restructuring, but actively delays it by manufacturing the appearance that no restructuring is needed.

\subsection*{A brief illustration outside infrastructure}

This concept is given a single, deliberately modest illustration from a third domain, offered as exactly that --- one additional anecdotal case extending the concept's apparent reach, not a third leg of the formal two-domain comparison Chapter~\ref{ch:two-domains} builds and is careful to keep at $n=2$. Shannon Vallor, Baillie Gifford Chair in the Ethics of Data and Artificial Intelligence at the Edinburgh Futures Institute, speaking at the ``Is AI a Threat?'' panel hosted by the Royal Institute of Philosophy on 26 June 2026, describes a finding from two independent surveys of executives (roughly two hundred UK chief executives in one study, three hundred US chief executives in another) using AI systems for business decisions: the large majority used AI for most of their decisions, trusted its advice over that of their human peers, and --- in about half of cases --- deferred to the AI's conclusion over their own judgment when the two conflicted. Vallor's own term for the end state of this pattern is \emph{cognitive surrender}: a condition in which a decision-maker believes themselves to still be exercising judgment while the substantive work of recognition and implementation has, in practice, migrated to the system they are nominally only consulting.

The structural parallel to this chapter's central distinction is direct enough to be worth stating, while resisting the temptation to claim more than the single case supports. ``Keeping a human in the loop'' is, on this account, a nominal closure: $N(t) = 1$, an oversight mechanism is declared present. Whether $O(t) = 1$ --- whether the human's presence is doing any actual epistemic work, rather than rubber-stamping a conclusion already reached elsewhere --- is a separate, empirical question the declaration does not answer, and Vallor's CEO data suggests that for a substantial fraction of a real, surveyed population, it does not. Asked directly whether keeping a human in the loop addresses this risk, Vallor's answer was that the phrase functions, at best, as what she called a spongy guardrail, and at worst a total illusion, precisely because cognitive surrender is the condition of believing oneself still in the loop while no longer being so --- a closure illusion stated in almost exactly this chapter's vocabulary, by a source that had no reason to be aware of it.

\section{What this chapter actually establishes}

Not the clean ordering it set out to test. What it establishes instead: tightly coupled subsystems can reach joint closure with negligible lag, consistent with --- though not independently confirming --- the coupling--lag conjecture; loosely coupled or independently-originated subsystems may share no common reference point at all, meaning lag is sometimes simply the wrong concept to reach for; and declared closure and operational closure are separate quantities that can diverge for extended periods, sometimes because a local patch actively obscures the divergence. These three findings recur, in a different costume, in the chapter that follows.
